\documentclass[compacto,10pt]{pubciencias-notas}

% -- Paquetes adicionales
\usepackage{comandosEPN}

% -- Datos del libro
\autor{Departamento de Formación Básica}
\materia{Álgebra Lineal}

\nota{Bases del Concurso}
\fecha{Semestre 2024-B}

\logouno[1.1cm]{Logo_EPN.eps}
\logodos[1.7cm]{Logo_DFB.eps}

\definecolor{colordef}{cmyk}{0.81,0.62,0.00,0.22}

\begin{document}

\encabezado

\vspace*{-4mm}

\setcounter{section}{3}
\section{CONCURSO DE ÁLGEBRA LINEAL}
\setcounter{subsection}{5}
\subsection{Aspectos Generales}

El \textbf{concurso de Álgebra Lineal} es un concurso de retos y trivias estudiantiles cuyo objetivo es fomentar el conocimiento de Álgebra Lineal de una manera divertida, estimulando a que los estudiantes desarrollen sus destrezas cognitivas en el ámbito de la ciencia y la tecnología.


Para participar en el concurso se deberá:
\begin{enumerate}

\item 
Ser estudiante universitario.
\item 
Contar con un grupo de hasta (4) estudiantes.
\item 
Contar con un celular por grupo, de preferencia que tenga datos.
\item 
Llenar el formulario de inscripción con su grupo de hasta cuatro (4) estudiantes. 

\end{enumerate}
\end{document}